\section{Klasyfikacja metod obrony}

Wśród technik mających na celu zwiększenie bezpieczeństwa dużych modeli językowych wobec bezpośrednich ataków promptowych, można wyróżnić interwencje realizowane na etapie trenowania, wdrożenia w realne środowisko oraz jego ewaluacji~\cite{correia2026systematicliteraturereviewllm}. Podział ten stanowi rozszerzenie klasyfikacji przygotowanej przez agencję \gls{NIST}~\cite{nist_taxonomy_2025} i jest oparty na na przeglądzie literatury opracowanym przez Correi~i~in.~\cite{correia2026systematicliteraturereviewllm}. Ze względu na ograniczenie zakresu pracy do bezpośrednich ataków, skupiono się głównie na mitygacjach stosowanych na etapie wdrożenia i pominięto kategorie związane z pośrednimi atakami na duże modele językowe.

\subsection{Metody stosowane na etapie trenowania modelu}
Metody stosowane na etapie trenowania mają na celu zwiększenie odporności samego modelu. W tym przypadku wyodrębnić można strategie użyte podczas uczenia wstępnego, które jednak nie są często wykorzystywane ze względu na zbyt duży koszt~\cite{correia2026systematicliteraturereviewllm}. Wymienić można również strategie, których celem jest wyrównanie bezpieczeństwa (ang. safety alignment) wykonane już po uczeniu.
Wykorzystywana jest metoda \gls{SFT}, za pomocą którego model może być trenowany przy użyciu zbioru zawierające zapytania niebezpieczne wraz z oczekiwanymi odmowami. Stosowane jest także podejście oparte na~\gls{RLHF}, czyli uczeniu przez wzmacnianie korzystające ze sprzężenia zwrotnego od człowieka, które ma na celu zoptymalizowanie modelu względem ludzkich preferencji dotyczących bezpieczeństwa~\cite{ouyang2022traininglanguagemodelsfollow}. Aplikowana jest również technika \gls{DPO}, która pozwala na bezpośrednią optymalizację docelowego modelu na podstawie par preferencji ludzkich bez potrzeby stosowania osobnego modelu nagrody~\cite{rafailov2024directpreferenceoptimizationlanguage}. Istnieje także możliwość wyeliminowania sprzężenia zwrotnego od człowieka na rzecz informacji pochodzącej od sztucznej inteligencji, korzystając z~\gls{RLAIF}. Model nagrody trenowany jest w takim podejściu na podstawie etykiet generowanych przez inny model językowy, co eliminuje potrzebę czasochłonnej adnotacji danych przez człowieka~\cite{bai2022constitutionalaiharmlessnessai}.
Metody omówione w tej sekcji nie będą rozpatrywane w części eksperymentalnej pracy ze względu na ich wysoki koszt obliczeniowy.

\subsection{Metody stosowane na etapie wdrożenia modelu}
Metody stosowane na etapie wdrożenia modelu obejmują natomiast mechanizmy mające na celu przeciwdziałanie atakom, które ominąć mogą zabezpieczenia wprowadzone podczas treningu~\cite{correia2026systematicliteraturereviewllm}.
Ze względu na relatywnie niski koszt wdrożenia i brak konieczności trenowania modelu, metody należące do tej kategorii stanowią główny przedmiot analizy.

W tym celu stosowane są między innymi strategie oparte na przygotowaniu promptu, na podstawie których ukierunkować można zachowanie modelu i określić jego ograniczenia przy użyciu instrukcji systemowej lub jawnie oddzielać instrukcję od danych użytkownika czy przypominać o zasadach po poleceniu użytkownika.
Przykładem jest metoda Self-Reminder~\cite{wu2023selfreminder}, która oparta jest na dodaniu instrukcji przed i po poleceniu użytkownika, która przypomina modelowi, by odpowiadał zgodnie z zasadami bezpieczeństwa, co dla badanych przypadków obniżyło sukces ataków o około 48\%.

Stosowane również są strategie, w których przekształcane jest zapytanie użytkownika. Po wykonanych zmianach jest ono następnie przekazywane do modelu docelowego~\cite{correia2026systematicliteraturereviewllm}.
Polecenie modyfikowane jest na poziomie znaków poprzez ich losowe usunięcie lub podmianę. Dzielone mogą być także tokeny, co ma osłabić ataki oparte na ich konkretnych kombinacjach. Przekształcenia można również dokonać na poziomie zdania, a stosowane są w tym celu parafrazowanie czy tłumaczenie. Przykładem techniki wykorzystującej modyfikację wejścia jest metoda Backtranslation~\cite{wang2024defendingllmsjailbreakingattacks}, w której jeden model generuje odpowiedź wykorzystywaną przez drugi model do odtworzenia oryginalnego zapytania użytkownika. Dla modelu Llama-2-13b technika ta pomogła badaczom zwiększyć liczbę obronionych ataków GCG z 66\% do 100\% oraz ataków PAIR z 64\% do 94\%.

Następnym podejściem jest strategia oparta na analizie wielu odpowiedzi wygenerowanych na podstawie różnych wariantów zapytania użytkownika w celu wykrycia zagrożenia~\cite{correia2026systematicliteraturereviewllm}. Reprezentantem takiej techniki jest metoda SmoothLLM~\cite{robey2024smoothllmdefendinglargelanguage}. Przy jej użyciu w przypadku modelu Llama~2 zredukowano skuteczność ataku GCG do 0,1\%.

Istotną kategorią są metody obrony na poziomie samego modelu, które nie wymagają jego ponownego trenowania. Stosowane jest sterowanie procesem dekodowania (ang. decoding steering), w ramach którego modyfikowany jest w czasie rzeczywistym rozkład prawdopodobieństwa tokenów wyjściowych, by ukierunkować go na bezpieczniejszą odpowiedź~\cite{correia2026systematicliteraturereviewllm}.
Przykładem jest metoda SafeDecoding~\cite{xu2024safedecodingdefendingjailbreakattacks}, która łączy rozkłady prawdopodobieństwa modelu docelowego i drugiego modelu dostrojonego pod kątem bezpieczeństwa. Zredukowała ona dla modelu Llama~2 skuteczność ataku GCG z 32\% do 0\% oraz ataku PAIR z 18\% do 4\%.
Podejściem zaliczanym do tej kategorii są także metody wykorzystujące inżynierię aktywacji. W ramach metody InferAligner~\cite{wang2024inferalignerinferencetimealignmentharmlessness} wykorzystano wektory sterujące bezpieczeństwem (ang. safety steering vectors), które wydzielono z modelu dostrojonego pod kątem bezpieczeństwa w celu modyfikacji aktywacji modelu docelowego poczas inferencji. Dzięki temu badacze zmniejszyli skuteczność ataków na dostrojony model Llama~2 do blisko 0\%.
Wyróżnić można także techniki przycinania (ang. defensive pruning) oparte na kompresji modelu w celu usunięcia mało znaczących parametrów~\cite{correia2026systematicliteraturereviewllm}, co może skutkować zwiększeniem odporności modelu na ataki~\cite{hasan2024pruningprotectionincreasingjailbreak}.

Kolejną kategorią metod stosowanych na etapie wdrożenia modelu jest wykrywanie i przerywanie szkodliwych interakcji. W jej ramach wyszczególnić można filtrowanie wejścia i wyjścia. Polega na zastosowaniu zewnętrznego mechanizmu klasyfikującego zapytania lub odpowiedzi jako szkodliwe~\cite{correia2026systematicliteraturereviewllm}. Podejścia te wykorzystują zewnętrzne modele moderujące, które są dostrojone do sprawdzania bezpieczeństwa tekstu. Przykładem takiego modelu jest LlamaGuard~\cite{inan2023llamaguardllmbasedinputoutput}.
Do tej kategorii zalicza się również techniki oparte na analizie miary perplexity, dzięki którym można wykryć ataki oparte na dodawaniu wygenerowanych sufiksów do niebezpiecznego zapytania~\cite{alon2023detectinglanguagemodelattacks}.
Wymienić można także podejście oparte na metodzie samorefleksji modelu. Opiera się na technice łańcucha rozumowania (ang. chain-of-thought), w której model językowy ocenia wygenerowaną przez siebie odpowiedź pod kątem niedozwolonych treści. Podejście wykorzystujące tę technikę zredukowało sukces ataku na zbiorze ataków DAN z 48,05\% do 0,15\%~\cite{cao2024chainofthought}.

\subsection{Metody stosowane na etapie ewaluacji modelu}
Metody stosowane na etapie ewaluacji modelu służą przede wszystkim ocenie podatności modelu~\cite{correia2026systematicliteraturereviewllm}. Nie przeciwdziałają one bezpośrednio atakom, ale pomagają identyfikować słabe punkty obrony. Do tej kategorii zaliczają się między innymi zestawy ataków, zbiory szkodliwych zapytań czy porównania służące do oceny skuteczności obron. Jednymi z najczęściej wykorzystywanych ataków w pracach naukowych są ataki GCG, PAIR i AutoDAN~\cite{correia2026systematicliteraturereviewllm}. Używane są jako punkt odniesienia. Często stosowane są zbiory ataków do ewaluacji bezpieczeństwa. Są nimi AdvBench~\cite{zou2023universaltransferableadversarialattacks}, JailbreakBench~\cite{chao2024jailbreakbenchopenrobustnessbenchmark} oraz HarmBench~\cite{mazeika2024harmbenchstandardizedevaluationframework}.